\documentclass[11pt]{article}
\usepackage[margin=1in]{geometry}
\usepackage[T1]{fontenc}
\usepackage{lmodern}
\usepackage{amsmath}
\usepackage{hyperref}

\title{Journal Entry: Current Project Status}
\author{calc\_mbcwc}
\date{2026-05-20}

\begin{document}
\maketitle

\section*{Status}

The project now has a working first benchmark target for the Stanford matrix
\(\Phi^4\) Schwinger-Keldysh calculation. The core result is implemented in
\texttt{scripts/stanford\_phi4\_kernel.py} and tested in
\texttt{tests/test\_stanford\_phi4\_kernel.py}.

The benchmark is model driven. It starts from the Stanford matrix \(\Phi^4\)
model, applies the standard Schwinger-Keldysh branch lift
\[
  \phi_1 = \phi_r + \frac{1}{2}\phi_a,
  \qquad
  \phi_2 = \phi_r - \frac{1}{2}\phi_a,
\]
derives the normalized branch-difference coefficients of
\[
  \phi_1^4 - \phi_2^4,
\]
records the on-shell ladder-kernel equations from the Stanford paper, and
evaluates the leading fixed-mass and thermal-mass-improved massless behavior.

\section*{Recovered Stanford Answers}

The current pure-Python benchmark reproduces
\begin{align}
  m_{\mathrm{th}}^2 &= \frac{2\lambda}{3\beta^2},\\
  \lambda_L &\simeq 0.025\,\frac{\lambda^2}{\beta^2 m},
  \qquad m\beta \ll 1,\\
  \lambda_L &\simeq 0.031\,\frac{\lambda^{3/2}}{\beta},
  \qquad m = 0,
\end{align}
where the massless expression uses the thermal-mass improvement.

The Appendix kernel diagonalization also reproduces the small-mass coefficient
numerically. For the tested \(m\beta=0.1\) grid,
\[
  \frac{\beta^2 m}{\lambda^2}\lambda_L = 0.0256986096,
\]
consistent with the quoted Stanford value near \(0.025\).

\section*{Math Audit}

The math audit found and corrected important numerical details. The diagonal
\(E_-=0\) case in the Appendix kernel now uses the analytic limit rather than
an artificial epsilon. The inverse hyperbolic-sine factors are evaluated
stably at large argument. The default momentum mapping scale \(P_0\) is now
mass-regime aware, using the Stanford recommendation \(P_0 \approx 3m\) in the
small-mass regime and \(P_0 \approx m\) otherwise. The numerical matrix is
checked for symmetry before using a symmetric eigensolver.

The symbolic audit clarified that the quantities previously called SK vertex
weights are normalized branch-difference coefficients only. They do not include
the interaction sign, coupling, or real-time factor required for full Feynman
vertex factors. The code now exposes that distinction in the pipeline payload.

\section*{Verification}

The current verification evidence is:
\begin{itemize}
  \item \texttt{python3 -m unittest tests/test\_stanford\_phi4\_kernel.py}:
  20 tests pass.
  \item A scoped non-Wolfram regression suite: 51 tests pass.
  \item Full discovery still has known \texttt{WolframKernel}
  metadata-export segfaults in the Wolfram-dependent tests. Those failures are
  separate from the pure-Python Stanford benchmark and remain a runtime
  stability issue to solve later.
\end{itemize}

\section*{Next Target}

The next research target is to connect the benchmark pipeline to a general
model-to-Schwinger-Keldysh-kernel pipeline. The current Stanford implementation
is an acceptance reference: future generalization should keep reproducing these
answers automatically while replacing Stanford-specific assumptions with data
derived from an input model.

\end{document}
